\input BaTeX.tex

\PapierGroesse(66pt,100pt)

\modul{ean,pstex}
\parskip=0pt
\begin{document}
\psSetze{linienstil=}
\begin{psfigur}(64,98)
\psPin(0,0){
\psCode{0.1 setflat}
\psPolygon[fuellfarbe=0.3,eckradius=4,linienstil=](6,1)(58,1)(64,50)(64,85)(62,98)(2,98)(0,85)(0,50)
\psPolygon[fuellfarbe=0.7:0.9:0.8,linienstil=](6,52)(58,52)(58,79)(6,79)
\psMultipin(15,14)(12,0){3}{%
	\psMultipin(0,0)(0,9){4}{%
		\psPolygon[fuellfarbe=0.8,eckradius=1.5,linienstil=](0,0)(10,0)(10,7)(0,7)
	}
}
\psPin(15,5){
	\psPolygon[fuellfarbe=0.9:0.2:0.2,eckradius=1.5,linienstil=](0,0)(10,0)(10,7)(0,7)
}
\psPin(27,5){
	\psPolygon[fuellfarbe=0.9:0.9:0.2,eckradius=1.5,linienstil=](0,0)(10,0)(10,7)(0,7)
}
\psPin(39,5){
	\psPolygon[fuellfarbe=0.2:0.9:0.2,eckradius=1.5,linienstil=](0,0)(10,0)(10,7)(0,7)
}
\psPin(30,88){
	\psSetze{linienrand=c,farbe=0.9,linienbreite=2,linienstil=-}
	\psKreisbogen[farbe=0.9,linienbreite=1.25,linienstil=-](0,0){1.5}(-45,45)
	\psKreisbogen[farbe=0.9,linienbreite=1.25,linienstil=-](0,0){4}(-45,45)
	\psKreisbogen[farbe=0.9,linienbreite=1.25,linienstil=-](0,0){6.5}(-45,45)
	\psKreisbogen[farbe=0.9,linienbreite=1.25,linienstil=-](0,0){9}(-45,45)
}
}
\end{psfigur}


\end{document}